\documentclass[11pt]{article}
\usepackage[margin=1in]{geometry}
\usepackage{amsmath,amssymb}
\usepackage{hyperref}
\usepackage{listings}
\usepackage{xcolor}
\usepackage{parskip}

\lstset{
  basicstyle=\ttfamily\small,
  keywordstyle=\color{blue},
  commentstyle=\color{gray},
  frame=single,
  breaklines=true,
}

\title{Design Document: Align Mode for WhyRel}
\author{WhyRel Project}
\date{\today}

\begin{document}
\maketitle

\tableofcontents
\newpage

% ============================================================
\section{Overview}
% ============================================================

Align mode (\texttt{-align}) is a tool-assisted workflow within WhyRel for constructing
\emph{bicommands} --- the relational programs that pair up the execution steps of two
unary methods and serve as the proof skeleton for relational verification.

Writing bicommands by hand is tedious and error-prone.  The alignment of two programs
is not unique: many correct alignments exist, and the ``best'' one for a given proof
depends on the relationship being proved.  Align mode addresses this by:

\begin{enumerate}
  \item Computing a \emph{default alignment} automatically from the two method bodies.
  \item Displaying the result in a form the user can inspect and refine.
  \item (Planned) Providing an interactive rewriting interface so the user can pivot to
        a better alignment without leaving the tool.
\end{enumerate}

% ============================================================
\section{Background}
% ============================================================

\subsection{Relational verification in WhyRel}

WhyRel verifies relational properties (equivalence, noninterference, simulation) by
asking the user to supply a \emph{bimodule}: a pair of modules together with a
relational specification and a \emph{bicommand} body.  A bicommand is a structured
interleaving of two unary command sequences; it must project onto the respective
unary bodies when the left or right side is erased.

The two primitive bicommand constructors are:
\begin{description}
  \item[\texttt{|[c, d]|}] (\texttt{Bisync}/\texttt{Biprod}) --- execute $c$ and $d$
        synchronously; both sides advance together.
  \item[\texttt{|[c, skip]|} / \texttt{|[skip, d]|}] (\texttt{Bisplit}) --- advance
        only one side; the other idles.
\end{description}

Higher-level constructs (\texttt{Bivardecl}, \texttt{Biwhile}, \ldots) handle
variable scoping and loops.

\subsection{The alignment problem}

Given unary method bodies $L$ and $R$, an alignment is any bicommand $B$ such that
$\pi_L(B) \equiv L$ and $\pi_R(B) \equiv R$, where $\pi_L$, $\pi_R$ are the left and
right projection functions.  Alignments that pair corresponding steps with
\texttt{Bisync} tend to produce the easiest proofs; purely split alignments always
work but give weak relational structure.

% ============================================================
\section{Current Implementation}
% ============================================================

\subsection{Entry point}

Align mode is invoked with:
\begin{lstlisting}
whyrel -align -l <LMod> <LMeth> -r <RMod> <RMeth> files... [-o <file>] [-man]
\end{lstlisting}

\texttt{main.ml} parses and type-checks the input files in the usual way, then calls
\texttt{Align.run} instead of the translation pipeline.

\subsection{Default alignment (\texttt{align.ml})}

The default alignment is computed by \texttt{compose\_sequentially}:

\begin{enumerate}
  \item Look up the bodies of \texttt{LMod.LMeth} and \texttt{RMod.RMeth} in
        \texttt{penv} via \texttt{find\_method\_body}.
  \item Call \texttt{build\_default lc rc}, which:
        \begin{itemize}
          \item Pierces matching \texttt{Vardecl} nodes on either side with
                \texttt{Bivardecl}, preserving scoping.
          \item Flattens remaining sequences with \texttt{flatten\_seq}.
          \item Interleaves the flat step lists with \texttt{interleave\_steps}: each
                left step is emitted as \texttt{|[c, skip]|} followed by the
                corresponding right step as \texttt{|[skip, d]|}; unmatched tail
                steps are padded with \texttt{skip}.
        \end{itemize}
\end{enumerate}

This produces a fully split alignment --- no synchronisation is assumed.  It is always
valid (projections are correct by construction) but typically not the intended proof
structure.

\subsection{Output modes}

Two output modes are currently supported, selected by the \texttt{-man} flag.

\paragraph{Default mode (HTTP server, \texttt{http\_server.ml}).}
The bicommand is pretty-printed to a string and served as plain text over a local HTTP
server on port 8080 at \texttt{/bicom}.  The client (e.g.\ a browser or \texttt{curl})
can fetch the alignment text.

\paragraph{Manual / HTML mode (\texttt{-man}, \texttt{html\_server.ml}).}
The bicommand is rendered as a line-numbered HTML page (dark-theme monospace) and
served at \texttt{/bicom}.  This is useful for visual inspection.

% ============================================================
\section{Planned Features}
% ============================================================

\subsection{Synchronisation heuristics}

The default split-everything alignment is rarely the desired proof structure.  Common
heuristics to improve it:

\begin{itemize}
  \item \textbf{Syntactic matching}: pair steps that are syntactically identical
        (same assignment shape, same field, same operator) as \texttt{Bisync}.
  \item \textbf{Loop pairing}: when both sides contain a \texttt{While} loop at the
        same position, emit a \texttt{Biwhile} rather than two split loops.
  \item \textbf{User-provided hints}: allow the user to annotate the source with
        alignment hints that override the default.
\end{itemize}

\subsection{Rewriting rules (\texttt{rewrites.ml})}

Once a default alignment is displayed, the user should be able to refine it by applying
named rewriting rules that preserve projection correctness.  Candidate rules include:

\begin{description}
  \item[\textsc{Sync}] Merge \texttt{|[c, skip]| ; |[skip, d]|} into \texttt{|[c, d]|}
        when $c$ and $d$ are ``related enough'' (e.g.\ same shape).
  \item[\textsc{SplitL} / \textsc{SplitR}] Reverse of \textsc{Sync}: split a
        \texttt{Bisync} step into two split steps.
  \item[\textsc{SwapL} / \textsc{SwapR}] Move a split step past the next step on the
        same side (commuting independent statements).
  \item[\textsc{PairWhile}] Convert two consecutive \texttt{While} bisplit loops into a
        \texttt{Biwhile}.
  \item[\textsc{UnpairWhile}] Reverse of \textsc{PairWhile}.
\end{description}

Each rule has a \emph{focus}: the line or node in the bicommand where it applies.
Rules should be invertible where possible to support undo.

\subsection{Interactive mode (\texttt{interactive.ml})}

The interactive mode extends the HTTP server with a stateful session:

\begin{enumerate}
  \item The server holds the current bicommand as mutable state.
  \item \texttt{GET /bicom} returns the current alignment (HTML or plain text).
  \item \texttt{POST /rewrite} applies a named rewrite rule at a given focus; the
        server returns the updated alignment.
  \item \texttt{POST /undo} reverts the last rewrite.
  \item \texttt{GET /export} writes the current bicommand as a WhyRel bimodule
        skeleton to the output file.
\end{enumerate}

A thin browser-side UI (or a CLI companion) can drive this API, letting the user
iteratively refine the alignment before copying the result into their \texttt{.rl}
file.

\subsection{Export to bimodule skeleton}

After alignment is finalised, align mode should emit a \texttt{.rl} file containing:
\begin{itemize}
  \item The biinterface derived from the two module interfaces.
  \item A bimodule stub with the aligned bicommand body and placeholder relational
        specifications (\texttt{requires true}, \texttt{ensures true}).
\end{itemize}
This skeleton can be opened in an editor, filled in with the actual relational
specifications, and fed back to the full WhyRel pipeline.

% ============================================================
\section{Architecture}
% ============================================================

\begin{verbatim}
Align.run
  |
  +-- compose_sequentially     (align.ml)
  |     |
  |     +-- find_method_body   lookup in penv
  |     +-- build_default      structural recursion over Vardecl/Seq
  |           +-- interleave_steps   flat step pairing
  |
  +-- [man_mode = false]  start_server (http_server.ml)   plain text
  +-- [man_mode = true]   start_html_server (html_server.ml) HTML view
  |
  (planned)
  +-- Interactive.session      stateful rewrite loop
  |     +-- Rewrites.apply     rewrite rules
  |     +-- Rewrites.undo      history stack
  +-- Export.bimodule_skeleton emit .rl stub
\end{verbatim}

% ============================================================
\section{Data Flow}
% ============================================================

\begin{enumerate}
  \item \textbf{Input}: \texttt{penv} (type-checked program environment),
        left/right module and method names.
  \item \textbf{Body extraction}: \texttt{find\_method\_body} returns
        \texttt{Annot.command option} for each side.
  \item \textbf{Alignment}: \texttt{build\_default} produces an
        \texttt{Annot.bicommand}.
  \item \textbf{Serialisation}: \texttt{Pretty.pp\_bicommand} renders the
        bicommand to a string.
  \item \textbf{Serving}: the string (or HTML-wrapped version) is served over
        HTTP.
  \item \textbf{(Future) Editing}: rewrite rules transform the bicommand
        in-place; the updated string is re-served.
  \item \textbf{(Future) Export}: the final bicommand is written back to a
        \texttt{.rl} file.
\end{enumerate}

% ============================================================
\section{Design Decisions and Open Questions}
% ============================================================

\begin{enumerate}
  \item \textbf{Projection correctness as invariant.}  Every rewrite rule must
        preserve the projection invariant ($\pi_L(B) \equiv L$, $\pi_R(B) \equiv R$).
        Rules should be implemented with this as a postcondition and ideally checked
        by running the WhyRel projection checker after each rewrite.

  \item \textbf{Granularity of focus.}  Should rewrites target individual
        \texttt{Bisplit} nodes (fine-grained) or top-level sequence positions
        (coarse-grained)?  Fine-grained targeting is more expressive but requires
        a richer path/zipper representation.

  \item \textbf{State representation.}  The interactive server holds mutable OCaml
        state.  Should this be a simple reference to the current bicommand plus a
        history stack, or something more structured (e.g.\ a zipper for focus
        tracking)?

  \item \textbf{Conflict with the normal pipeline.}  Align mode currently skips
        pre-translation and translation.  If the user wants to verify the resulting
        bimodule, they must run WhyRel again on the exported \texttt{.rl} file.
        A tighter integration (e.g.\ align $\to$ verify in one invocation) could
        be explored.

  \item \textbf{Loop alignment.}  Biwhile pairing requires both sides to have loops
        at the same structural position.  What should happen when loop bounds differ
        or loop bodies have different shapes?

  \item \textbf{Port conflicts.}  Both server modes hardcode port 8080.  This should
        be configurable via a \texttt{-port} flag.
\end{enumerate}

\end{document}
